\chapter*{Resumen}

\section*{\tituloPortadaVal}

La lectura de ondas cerebrañes mediante electroencefalografía (EEG) lleva años existiendo, recientemente su integración con modelos de inteligencia artificial (IA) ha abierto nuevas vías para interpretarlas.

Hasta la fecha, los sistemas de interfaz cerebro-computador (BCI) han obtenido resultados fiables mediante paradigmas reactivos que dependen de estímulos externos constantes, lo cual no es aplicable a un sistema dinámico como puede ser un videojuego. 

Son pocos los estudios que han obtenido resultados sin estímulos externos (mediante imaginación motora, MI) y, aquellos que lo consiguen, hacen uso de hardware poco accesible por su coste o portabilidad.

En este trabajo se plantea como objetivo el diseño e implementación de un sistema BCI modular capaz de procesar señales EEG en tiempo real y traducirlas en comandos para controlar videojuegos. El sistema emplea un electroencefalógrafo ligero y asequible junto con el modelo de aprendizaje profundo MiRepNet, adaptado mediante técnicas de \textit{fine-tuning} para capturar las particularidades de cada usuario. 

Los resultados experimentales muestran una precisión media del 80\% en la clasificación binaria de imaginación motora (mano derecha frente a mano izquierda) sobre datos grabados (\textit{offline}), alcanzando valores superiores al 90\% en sujetos con alta calidad de señal. 

Para validar el sistema en un contexto dinámico, se han desarrollado dos videojuegos (\textit{Endless Wave Runner} y \textit{Project Arrows}) que integran filtros de decisión y mecanismos de asistencia o \textit{``aimbots''}. Estos mecanismos permiten compensar la variabilidad de las señales cerebrales, logrando una experiencia de usuario fluida y demostrando que es posible el control mental de aplicaciones complejas con hardware accesible.

\section*{Palabras clave}
   
\noindent BCI, imaginación motora (MI), inteligencia artificial, EEG, videojuegos, tiempo real, procesamiento de señales, accesibilidad, MiRepNet, interacción humano-computador.
